\section{Introduction}\label{sec:introduction}

Let $\mathcal P$ denote the set of positive primes, and put
$\N_0=\{0,1,2,\ldots\}$.  For $A,B\subseteq\N_0$, write
$A+B=\{a+b:a\in A,\ b\in B\}$.  Two subsets of $\N_0$ are
\emph{asymptotically equal} if their symmetric difference is finite.
Ostmann's conjecture asserts that $\mathcal P$ is asymptotically
additively indecomposable: it is not asymptotically equal to $A+B$ when
both summands contain at least two elements.  The conjecture goes back
to his 1956 treatise~\cite[p.~13]{Ostmann1956}; the eventual-equality
formulation and its terminology are recorded explicitly in
\cite[Definition~1.1 and Conjecture~1.2]{ElsholtzHarper}.
It is also known as the inverse Goldbach problem.

\begin{theorem}[Ostmann's conjecture]\label{cor:ostmann}
If $A,B\subseteq\N_0$ satisfy $|A|,|B|\ge2$, then
$(A+B)\mathbin{\triangle}\mathcal P$ is infinite.
Equivalently, every finite modification of $\mathcal P$ is additively
indecomposable into two sets with at least two elements each.
\end{theorem}

